\section{Provisional audit of the inference runtime}
\label{sec:audit}

\subsection{Purpose and method}

The literature motivates a memory-placement experiment, but it does
not determine whether \lcpp{} must be modified.  This audit asks a more
limited question: which controls are present in the source snapshot,
and which experiments can be performed without a patch?

The subject is upstream \lcpp{} commit
\path{ff067f76dd8e9e05f0528056f1274adf01a54d70}, inspected on
26 July 2026~\cite{llamacpp2026}.  The audit covers CMake
configuration, CPU buffer types, model loading, the KV-cache allocator
and NUMA argument handling.  File and line references are pinned to that
commit.  They describe source structure, not observed behaviour on
CRESCO8; all material findings require a build and residency check
before they are quoted as thesis results.

\subsection{What is already sufficient for the primary experiment}

A controlled whole-process comparison does not require native
memory-tier objects in the runtime.  With Flat mode verified, an external
launcher can bind both threads and memory to one socket's HBM node or
the corresponding DDR node.  A controlled page state then permits the
same binary, model, quantization and command to be compared under the
two policies.

This approach has three advantages.  It changes no runtime code, keeps
the causal treatment small, and establishes whether a tier effect is
large enough to justify finer placement work.  Its main risks are
file-backed model pages that pre-exist on another node, first-touch by
an unintended thread, automatic NUMA balancing, and hidden packed
copies.  These are measurement problems to verify with page-residency
tools, not reasons to begin with a runtime rewrite.

\subsection{Static findings}

\begin{table}[htbp]
\centering
\footnotesize
\caption{Static capability audit at commit \commitshort{}.  ``Not
exposed'' means no stable user control was found in the inspected path;
it does not mean that external Linux placement is impossible.}
\label{tab:runtime-audit}
\begin{tabularx}{\textwidth}{@{}P{2.8cm}P{5.35cm}Y@{}}
\toprule
\textbf{Capability} & \textbf{Source evidence} &
\textbf{Consequence for the thesis} \\
\midrule
One CPU backend device
& \path{ggml/src/ggml-cpu/ggml-cpu.cpp:508--524} returns one device
  and accepts only index zero.
& The scheduler does not represent individual CPU NUMA domains as
  devices.  External binding remains available. \\
\addlinespace
HBM build option
& \path{ggml/CMakeLists.txt:151} defines
  \texttt{GGML\_CPU\_HBM}.  The macro is private to the
  \texttt{ggml-cpu} target
  (\path{ggml/src/ggml-cpu/CMakeLists.txt:88--95}), while the default
  allocator is compiled in \texttt{ggml-base}
  (\path{ggml/src/CMakeLists.txt:192--209}).
& A successful HBM-enabled build is not evidence that default model
  allocations use HBM.  Confirm with symbol/build inspection and
  page residency. \\
\addlinespace
Dedicated HBM buffer type
& \path{ggml/src/ggml-cpu/hbm.cpp:24--54} implements a
  \texttt{CPU\_HBM} allocator, but the CPU extra-buffer list at
  \path{ggml/src/ggml-cpu/ggml-cpu.cpp:42--73} does not register it.
& The source contains useful plumbing, but no reachable
  \texttt{CPU\_HBM} placement target was found in the ordinary model
  loading path. \\
\addlinespace
Per-tensor override
& \texttt{-ot}/\texttt{--override-tensor} maps a tensor-name pattern
  to a buffer type, but the user-visible map is populated from device
  default buffer types
  (\path{common/arg.cpp:249--281}).
& The mechanism can select devices such as CPU/GPU; it does not expose
  local CPU HBM versus local CPU DDR in this snapshot. \\
\addlinespace
KV-cache placement
& For CPU execution, each layer starts with
  \texttt{ggml\_backend\_cpu\_buffer\_type()}; an offloaded layer uses
  its device buffer
  (\path{src/llama-kv-cache.cpp:210--223}).
& KV data type can be configured, but a separate CPU HBM/DDR placement
  selector was not found. \\
\addlinespace
NUMA control
& \texttt{--numa} offers \texttt{distribute}, \texttt{isolate} and
  \texttt{numactl} modes
  (\path{common/arg.cpp:2551--2564}).  The implementation enumerates
  NUMA CPUs and applies thread affinity
  (\path{ggml/src/ggml-cpu/ggml-cpu.c:636--717,2148--2191}).
& Runtime NUMA support should not be confused with memory-tier
  placement.  Use the external process memory policy and verify it. \\
\addlinespace
AMX path
& The AMX buffer registers quantized matrix multiplication for selected
  formats, including Q4, Q5, Q6 and Q8 families under shape constraints
  (\path{ggml/src/ggml-cpu/amx/amx.cpp:19--40,143--190}).
& AMX support exists, but kernel coverage for the selected model and
  phase must be demonstrated from build features, logs, disassembly or
  counters. \\
\bottomrule
\end{tabularx}
\end{table}

\subsection{The HBM option requires runtime confirmation}

The most consequential source observation concerns
\texttt{GGML\_CPU\_HBM}.  The generic CPU buffer allocator calls
\texttt{ggml\_aligned\_malloc} in
\path{ggml/src/ggml-backend.cpp:2299--2313}.  That function contains an
\texttt{hbw\_posix\_memalign} branch in
\path{ggml/src/ggml.c:331--383}.  At the pinned commit, however,
\texttt{ggml.c} belongs to \texttt{ggml-base}, while the HBM macro is
defined privately on the separate CPU backend target.  Static reading
therefore predicts that the generic allocator is compiled without the
HBM branch.

There is also a dedicated HBM buffer type in \path{hbm.cpp}; it is
compiled when the CPU backend is built, but no registration in the
CPU's extra-buffer list was found.  Taken together, these facts predict
that enabling the CMake option can advertise \texttt{CPU\_HBM=1}
without changing the default model allocation path.  This should be
written as a \emph{prediction}, not as a measured defect, until two
builds are compared on CRESCO8.

A minimal confirmation consists of:

\begin{enumerate}
  \item retain CMake output and the final compile commands for
        \texttt{ggml.c}, \texttt{hbm.cpp} and the CPU backend;
  \item inspect the resulting symbols or disassembly for calls to
        \texttt{hbw\_posix\_memalign};
  \item run an identical small model with the option off and on;
  \item record per-node page residency and throughput under a policy
        designed to distinguish HBM from DDR; and
  \item repeat with memory mapping disabled if pre-existing file-backed
        pages prevent an interpretable first-touch result.
\end{enumerate}

If the prediction is wrong, the audit should be corrected rather than
defended.  A static audit is valuable precisely because it supplies a
falsifiable test.

\subsection{Independent weights and KV placement}

Selective tiering is more demanding than whole-process binding.
\lcpp{} already has two useful abstractions: model tensors may be
overridden by name pattern, and KV buffers are created per layer and per
device.  In the inspected command path, neither abstraction exposes two
CPU memory tiers.  Consequently, the user cannot state ``weights in
local HBM and KV in local DDR'' as a normal command-line policy.

This is a real capability limitation for the proposed capacity study,
but it is not yet a performance gap.  It becomes a thesis-relevant
software problem only if all of the following hold:

\begin{enumerate}
  \item the whole-process HBM baseline is bandwidth-limited and faster
        than DDR;
  \item the desired live set exceeds usable local HBM;
  \item existing operating-system placement cannot express a stable,
        auditable split; and
  \item a plausible selective policy can be evaluated against the best
        unmodified baseline.
\end{enumerate}

The smallest likely intervention is to register distinct CPU HBM and
DDR buffer types and make the selected type reachable by the existing
override machinery.  A separate KV selector may then be needed.
Representing every NUMA domain as a backend device is a much larger
architectural change and should not be assumed necessary before
measurement.

\subsection{AMX and the Aurora result}

Goto et al.\ show that oneMKL uses AMX in CPU-side HPL-MxP work
~\cite{goto2026aurorahpl}.  This does not imply that the chosen
\lcpp{} model follows the same kernels.  The pinned source has an AMX
path for selected quantized matrix multiplications, including repacked
weights and shape restrictions.  The experiment must therefore report
which model tensors use the AMX buffer, whether the path differs between
prefill and decode, and how much extra packed storage it creates.

This distinction matters for the memory-tier comparison.  A run that
silently changes kernel coverage or packed representation changes both
\(W\) and \(Q\) in the model of Section~\ref{sec:synthesis}.  Build
features, runtime logs and peak resident memory belong in every run
manifest.

\subsection{Audit conclusion}

The audit does not justify making a \lcpp{} fork the thesis objective.
It supports a staged conclusion:

\begin{itemize}
  \item the primary, whole-process local-HBM versus local-DDR comparison
        is feasible with an unmodified, pinned runtime and verified
        operating-system placement;
  \item native HBM support in the inspected snapshot is uncertain and
        should be tested, not assumed;
  \item independent CPU-side weight and KV placement is not exposed in
        the inspected user path; and
  \item a runtime contribution is warranted only if the baseline shows
        a material, reproducible opportunity that existing controls
        cannot express.
\end{itemize}
